\subsection{Annotation and Tool Extraction}
\label{sec:annotation-tool-extraction}

The corpus was annotated in multiple passes with evidence tracing as a hard
requirement. In the diagnostic pass, each benchmark was mapped to a primary
ethical dimension and assessed against the 14-rubric HEART codebook.
Annotators recorded source evidence from the paper, repository, dataset card,
or evaluation artifact, rather than relying on impressionistic judgments. In
the extraction pass, two primary experts annotated the reusable evaluation
mechanisms that each paper could contribute and assigned expert tool labels.
Other experts then spot-checked and adjudicated disagreements between the two
primary annotations. The expert team also co-authored the guidebook cards from
the benchmark evidence, so that later users can follow the same toolkit to
adapt benchmark items rather than only inspect our final labels.

The extracted mechanisms include, for example, ambiguous/disambiguated
contrast sets, safe/unsafe contrast pairs, contextual integrity schemas,
temporal factuality checks, workflow simulations, and multi-turn adversarial
protocols. These mechanisms were then mapped to the corresponding HEART repair
tools and preserved as benchmark-specific sub-tools in the machine-readable
codebook and website cards.

Inclusion and exclusion were performed through iterative expert screening and
consensus discussion rather than a one-shot independent coding protocol.
Accordingly, we do not report Cohen's \(\kappa\) or Krippendorff's \(\alpha\)
for corpus selection: the screening logs were designed to produce a curated
resource and repair taxonomy, not to estimate independent-coder reliability.
Disagreements were resolved by evidence-based adjudication, with claims
narrowed when artifacts were incomplete, non-reproducible, or only weakly
supported by the source. Quantitative expert consistency is instead reported
for the blinded mini-case quality assessment in Section~\ref{sec:expert-quality},
where experts compare original and HEART-guided revisions under the 14 rubrics.

\subsection{Dimensions, Rubrics, and Tools}
\label{sec:dimensions-rubrics-tools}

The HEART framework organizes benchmark evaluation and repair along three
axes. First, the five ethical dimensions are derived from international policy
and governance sources, including UNESCO principles, GDPR, and the EU AI Act,
rather than from benchmark naming conventions. The dimensions are
Human-Centeredness, Fairness and Inclusiveness, Safety and Reliability,
Trustworthiness and Controllability, and Privacy Protection. They correspond
to distinct evidence needs: fairness requires evidence about protected
interests and proxy attributes, privacy requires evidence about information
flows, authorization, purpose, and consent, and safety requires evidence about
hazard type, risk level, vulnerable subjects, and escalation boundaries.

Second, the 14 diagnostic rubrics are grouped into three layers that target
different failure modes. Content Validity asks whether the benchmark has a
valid ethical measurement target: R1 construct clarity defines what capability
or risk is measured; R2 normative grounding records the legal, policy,
professional, or community basis for judgment; R3 source provenance and
evidence fitness checks whether the data source is traceable and suitable; R4
context and stakeholder coverage asks whether actors, affected parties,
relationships, and purposes are represented; and R5 real-world harm linkage
connects the task to concrete harm pathways. Evaluation Design asks whether
the task can actually measure that construct: R6 scenario validity checks
realism and domain fit; R7 task-format fit checks whether multiple choice,
open generation, multi-turn interaction, or agent workflow matches the risk;
R8 ground-truth and disagreement design records answer keys, uncertainty, and
legitimate disagreement; R9 metric validity checks whether the metric measures
the intended failure mode rather than an easy proxy; and R10 evaluator
reliability checks whether human or model judges are calibrated and
reviewable. Governance Reliability asks whether the benchmark remains usable
and auditable over time: R11 data and annotation QA checks annotation,
filtering, and conflict-resolution procedures; R12 robustness against gaming
and contamination checks hidden sets, leakage, and leaderboard overfitting;
R13 documentation and reproducibility checks cards, scripts, licenses, and
parameter records; and R14 maintenance and update governance checks release
versioning, feedback channels, update triggers, and retirement rules.

Third, the toolkit includes 14 generic repair operations to keep the workbox
actionable and reusable across dimensions. T01 maps broad ethical claims to
measurable constructs and protected interests, such as reframing a privacy
item around patient confidentiality rather than only leaked strings. T02
records normative sources and decision boundaries, such as noting that a
clinically knowledgeable friend is not automatically an authorized recipient
of therapy notes. T03 fills scenario, stakeholder, and information-flow slots,
such as sender, recipient, data subject, purpose, authorization, and harm. T04
adds risk tiers and hazard categories, useful for crisis or self-harm advice
where escalation thresholds matter. T05 adapts cultural and cross-lingual
categories, such as replacing identity groups or institutional settings with
locally valid ones. T06 builds contrast sets, such as BBQ-style
ambiguous/disambiguated pairs or safe/unsafe prompt pairs. T07 builds coverage
matrices over capabilities, scenarios, stakeholders, and risk modifiers. T08
turns single-round tests into multi-turn or adversarial red-team protocols.
T09 simulates agents and workflows, such as a Notion search followed by a
Gmail send action. T10 verifies evidence and claims, such as checking whether
a time-sensitive CEO question has a stale false premise. T11 reports
multi-metric failure profiles instead of a single score, such as separating
privacy leakage, insufficient de-identification, and helpfulness. T12 uses
rubric-based and multi-judge scoring for open-ended outputs. T13 packages
benchmark cards, data cards, scripts, model versions, prompts, and reproduction
commands. T14 governs lifecycle, hidden splits, contamination checks, update
triggers, and retirement. Domain-specific mechanisms extracted from the corpus
are maintained as instantiated sub-tools rather than promoted into additional
generic tools, which keeps the core toolkit stable while preserving concrete
evidence from the 103 audited benchmarks.
